\section{Introduction}\label{sec:nintro}

Un topos peut se présenter de bien des manières qui semblent sans rapport: comme les actions d'un monoïde, comme les représentations d'un carquois, comme les faisceaux sur un espace, comme les modèles d'une théorie logique. Les \emph{ponts} d'O.~Caramello \cite{caramello-unification,caramello-tst} tirent parti de cette diversité. Une propriété du topos invariante par équivalence se caractérise sur chacune de deux présentations; ces caractérisations sont les \emph{arches} du pont, et l'équivalence entre les deux présentations en est le \emph{tablier}. Mises bout à bout, les deux arches donnent un théorème qui relie les deux présentations et où le topos n'apparaît plus.

Dans un exposé aux Mines Paris \cite{caramello-mines}, Caramello a proposé de confier à une machine la combinaison systématique des arches et des tabliers. L'article associé \cite{bridges-generator} fait tourner à la main une petite version de ce générateur et la décrit comme un pipeline: un registre d'arches, démontrées en Lean~4 sous la forme \texttt{(propriété de PSh C) ↔ (condition sur C)}; un registre de tabliers, petites catégories $C$ dont les préfaisceaux forment une structure familière; enfin des étapes qui associent chaque tablier à chaque arche, calculent la condition sur $C$, la réécrivent dans la langue du domaine et la vérifient en Lean, par une preuve directe et dans la littérature. La présente note applique ce générateur aux notes manuscrites de Grothendieck conservées à l'Université de Montpellier, où l'on trouve déjà des tabliers de ce genre et des invariants lus sur deux présentations d'un même topos.

\subsection*{Plan}

La section~\ref{sec:recall} rappelle les notations et les cinq arches de l'article associé. La section~\ref{sec:morearches} donne trois arches de plus (deux valeurs, connexe, spatial), démontrées en Lean, et la section~\ref{sec:grdecks} les lit sur trois tabliers qui reviennent dans les notes de Grothendieck. La section~\ref{sec:precedents} rassemble deux précédents tirés de ces notes: l'exercice du dossier 161-3, p.~26, que le générateur retrouve et dont nous démontrons en Lean la forme relative aux préordres, avec la même identification au dossier 37, p.~48; puis les sous-topos d'une théorie, au dossier 161-2. La section~\ref{sec:pseudofields} tire du dossier 161-2 un tablier nouveau, le topos des pseudo-corps, dont Grothendieck donne trois présentations; nous en vérifions une, nous y lisons les arches et nous répondons à une question qu'il a laissée en marge. La section~\ref{sec:stacks} relie les topos classifiants aux champs de Grothendieck et propose des arches plus fines, avec un exemple tiré du dossier 146. La section~\ref{sec:nconclusion} conclut, et l'annexe~\ref{sec:nlean} donne les fichiers et les déclarations Lean.
